\section{Research Taste en la era de la IA: lo que le queda al ser humano es el nivel de abstracción}

La parte intermedia de la conversación pasa de la deserción escolar a una pregunta más amplia: cuando la IA produce cada vez más artifacts, ¿cómo juzga el ser humano qué research es bueno? Zhao Chenyang (赵晨阳) y el conductor coinciden en que la applied research todavía puede reflejarse con revenue, valor para el usuario o benchmarks claros; la fundamental research es más difícil, porque se parece más a un arte y necesita personas con taste que la juzguen. En la era de la IA aumentan los artifacts, la capacidad de revisión humana es limitada y el peer review tradicional difícilmente resiste esa presión.

Aquí surge un juicio clave: el buen gusto científico se manifiesta en el nivel de abstracción con que se plantea un problema. Detrás de un problema concreto puede haber una estructura más fundamental, por ejemplo el consenso, la consistencia, la generalización o el aprendizaje a largo plazo dentro de la intelligence. La investigación que de verdad no es fácil de sustituir pronto por la IA no suele ser construir en seis meses un sistema soon-to-be-obsolete, sino lograr abstraer el problema de largo plazo que hay detrás de un testbed concreto.

\begin{importantbox}{El lugar del research taste}
Una vez que baja el costo de ejecución, el valor del ser humano no está en escribir más código o hacer más experimentos que la IA, sino en plantear abstracciones de problemas de más alto nivel, en juzgar qué problemas valen la pena, qué tradeoffs importan y qué rutas son solo optimizaciones locales.
\end{importantbox}

Zhao Chenyang comenta que le gusta escribir debug log y que mantiene un machine learning systems tutorial. Espera que en el futuro la IA pueda aprender de esos debug log el proceso de prueba y error. Este detalle es importante, porque los artículos académicos por lo general solo presentan la idea final ya pulida y rara vez registran los tradeoffs intermedios, los caminos fallidos y los juicios de depuración. Precisamente estos «juicios durante el proceso» quizá sean la razón por la que el taste es difícil de destilar y también difícil de sustituir con IA.

Esto se corresponde con el planteamiento de «la no linealidad de la investigación» de la entrevista de Xie Sanning (谢赛宁): cuando un artículo está escrito parece lineal, pero la idea real proviene de la exploración, el fracaso, los giros y la reinterpretación del problema. Si estos procesos no se registran, la IA solo puede aprender la forma del resultado y difícilmente puede aprender por qué el investigador abandonó cierta ruta, por qué conservó cierto diseño o por qué cierto fracaso en realidad aportaba mucha información.

\begin{knowledgebox}{Por qué importa el debug log}
\begin{itemize}
  \item Registra los caminos fallidos que un artículo no va a escribir.
  \item Expone los tradeoffs de diseño, en lugar de presentar solo la elección final.
  \item Convierte el «buen gusto», parcialmente, de una intuición indecible en material que se puede aprender.
  \item Puede llegar a ser una forma de datos importante para que la IA aprenda en el futuro el proceso de investigación.
\end{itemize}
\end{knowledgebox}

Este fragmento también señala una paradoja: una vez que la IA abarata la ejecución, un buen design se vuelve más importante. El costo de gobernar la escritura de código es mucho menor que el de contratar personas para ejecutar, siempre que el design sea suficientemente claro. La responsabilidad del ser humano ya no es escribir cada línea de código con sus propias manos, sino decidir la arquitectura, los límites, las prioridades, el manejo de fallos y los objetivos de largo plazo.

\subsection{Resumen de la sección}

En la era de la IA, el núcleo del research taste no es «escribir más», sino «abstraer más alto, juzgar con más precisión y registrar el proceso de manera más completa». La applied research puede apoyarse en la retroalimentación del mercado, mientras que la fundamental research depende más del taste y de una revisión de largo plazo. El debug log, el registro de tradeoffs y los caminos fallidos pueden llegar a ser material importante para que la IA futura aprenda el buen gusto en la investigación.
